\section{Actual zero-contact remainders under contact glueing}

Let $L=(N,E,I,d)$ be a marked ideal. Write $S(L)$ for the union of the
actual global components of $N$ contained in its actual cosupport, and
$R(L)=N\setminus S(L)$. A positive marking is essential in the following
zero-germ characterization. No bound on dimension, component count, or ideal
orders is imposed. Put
\[
 L^{\mathrm{rem},k}=(R(L),[\varnothing]^k\mathbin{+\!+}E_{\geq k},I,d).
\]

\begin{theorem}[The actual selected union is precisely the zero-germ locus]
\label{hir:zero-contact-selected-zero-germ}
\leanok
\lean{Hironaka.HironakaStepOneZeroContactRemainderGlueing.selectedComponents_iff_not_isNonzeroAlongAt,Hironaka.HironakaStepOneZeroContactRemainderGlueing.not_isNonzeroAlongAt_iff_vanishing_stalk}
\uses{hir:smooth-selected-components,hir:dimtwo-unselected-smooth-components}
Suppose $N$ is closed smooth, $I$ is coherent, and $d\geq1$. At each
$x\in N$, membership in $S(L)$ is equivalent to the failure of nonzero
along $N$ at $x$. Equivalently every germ of $I_x$ vanishes identically
on the germ $N_x$.
\end{theorem}
\begin{proof}
The selected and remaining component unions are disjoint closed smooth
sets with union $N$. If $x$ is selected, the complement of the remaining
union is a neighborhood of $x$ on which every point of $N$ lies in the
actual cosupport. A coherent ideal germ extends to a member at every
nearby point. At a nearby point of $N$, order at least $d\geq1$ makes
the value of this member zero. Thus no ideal member is frequently nonzero
along $N$ near $x$.

Conversely let $x$ be remaining. The actual cosupport is analytic, and
analytic continuation on each genuine connected smooth component proves
that its intersection with the remaining union contains no nonempty
relatively open subset. Near every remaining point, closedness of the
selected union makes the remaining and original carrier germs equal;
their along-carrier order predicates therefore agree. If all ideal germs
were zero along the remaining carrier at $x$, a finite coherent generating
family would give a single neighborhood where all ideal members vanish
along that carrier. This neighborhood would be contained in its cosupport,
contradicting the preceding thinness. A germ frequently nonzero along the
remaining carrier is also frequently nonzero along $N$.

Finally, the negation of frequent nonvanishing for every ideal member says
that each member is eventually zero along $N$. Coherence makes each member
holomorphic. These are exactly the two requirements for membership in the
vanishing-ideal stalk of $N$. Conversely such stalk membership contradicts
frequent nonvanishing. This proves the zero-germ equivalence.
\end{proof}

\begin{theorem}[Actual nonzero and selected germs commute with open restriction]
\label{hir:zero-contact-selected-open-dictionary}
\leanok
\lean{Hironaka.HironakaStepOneZeroContactRemainderGlueing.isNonzeroAlongAt_pullback_val_iff,Hironaka.HironakaStepOneZeroContactRemainderGlueing.remainingComponents_pullback_val}
\uses{hir:zero-contact-selected-zero-germ}
For a coherent ideal $I$, an arbitrary carrier $N$, an actual open $U$,
and $p\in U$, nonzero along the restricted carrier for $I|_U$ at $p$
is equivalent to nonzero along $N$ at the same ambient point. If $N$ is
closed smooth and $d\geq1$, the remaining carrier of $L|_U$ is precisely
the restriction of $R(L)$ to $U$.
\end{theorem}
\begin{proof}
Restrict an actual nonzero germ to $U$. Coherent ideal restriction preserves
its membership, and the open embedding identifies the neighborhood filter
at $p$ with that at its ambient point. It therefore preserves frequent
nonvanishing along the restricted carrier. In the reverse direction,
extend a restricted function back to the ambient manifold. Coherent
restriction identifies its ideal membership, while its values on $U$
are unchanged. The same equality of neighborhood filters transfers
frequent nonvanishing back. Apply the preceding zero-germ characterization
to the genuinely closed smooth restricted carrier and to $N$. Since
remaining membership means carrier membership and nonzero along, this
gives equality of the entire remaining carriers. No assertion about
equality of the two global component index sets is needed.
\end{proof}

\begin{theorem}[Actual nonzero and remaining germs commute with biholomorphisms]
\label{hir:zero-contact-selected-biholomorphic-dictionary}
\leanok
\lean{Hironaka.HironakaStepOneZeroContactRemainderGlueing.isNonzeroAlongAt_pullback_diffeo_iff,Hironaka.HironakaStepOneZeroContactRemainderGlueing.remainingComponents_pullback_diffeo}
\uses{hir:zero-contact-selected-zero-germ}
For a biholomorphism $\Phi$, coherent $I$, arbitrary $N$, and $p$ in its
source, nonzero along $\Phi^{-1}N$ for the actual entire pullback ideal
at $p$ is equivalent to nonzero along $N$ at $\Phi p$. For closed smooth
$N$ and positive marking, the remaining carrier of $\Phi^*L$ is
$\Phi^{-1}R(L)$.
\end{theorem}
\begin{proof}
Compose an actual ideal member with $\Phi$; it belongs to the entire
coherent pullback ideal. The homeomorphism's equality of neighborhood
filters transfers frequent nonvanishing. For the reverse implication
apply the same argument to $\Phi^{-1}$. Genuine coherent pullback
functoriality and the inverse identities identify the twice-pulled ideal
with the original ideal and the twice-pulled carrier with $N$. Closed
smoothness and coherence are preserved by this actual biholomorphism.
The zero-germ characterization on both sides now identifies remaining
membership at every point and hence the entire remaining carriers.
\end{proof}

\begin{theorem}[The entire remainder record commutes with actual transport]
\label{hir:zero-contact-remainder-transport-record}
\leanok
\lean{Hironaka.HironakaStepOneZeroContactRemainderGlueing.contactRemainderMarked_pullback_val,Hironaka.HironakaStepOneZeroContactRemainderGlueing.contactRemainderMarked_pullback_diffeo}
\uses{hir:zero-contact-selected-open-dictionary,hir:zero-contact-selected-biholomorphic-dictionary,hir:zero-contact-actual-remainder-definition}
For closed smooth $N$, coherent $I$, positive marking and every natural
$k$, the entire record $L^{\mathrm{rem},k}$ commutes with genuine open
restriction and with genuine biholomorphic pullback.
\end{theorem}
\begin{proof}
The preceding theorems give the exact carrier identities. The ideal and
mark are unchanged by the remainder construction, so their identities
follow from the definitions of actual pullback. On the boundary, taking
preimages commutes with dropping the first $k$ entries and with list
concatenation. The preimage of an empty set is empty, so the first $k$
historical entries remain exactly $k$ empty entries. These four field
identities give equality of the entire records. In particular a genuine
whole contact-record transport induces whole remainder-record transport.
\end{proof}

\begin{theorem}[True contact lifts transport the actual remainders at every prefix]
\label{hir:zero-contact-remainder-actual-contact-lifts}
\leanok
\lean{Hironaka.HironakaStepOneZeroContactRemainderGlueing.exists_contact_remainder_glueing}
\uses{hir:zero-contact-remainder-transport-record,mc:glueing-data-shrink,mc:contact-independence-along}
Suppose $K.N$ is closed smooth, $K.I$ is coherent, $K.d\geq1$, and
$x$ is in its actual cosupport. Let $u,v$ be genuine tangent directions
on actual opens containing $x$. There is actual glueing data $G$, with
both neighborhoods inside both tangent opens and containing $x$, such
that every actual test prefix $T$ of the coefficient-homogenized datum
has an actual lifted biholomorphism transporting the entire coefficient
contact remainder for $u$ to that for $v$, retaining every actual
exceptional entry and adding any prescribed number of empty historical
entries. This lift fixes the actual transformed cosupport. The padding
does not delete exceptional entries.
\end{theorem}
\begin{proof}
The actual homogenized contact construction gives glueing maps trivial
modulo the genuine glue ideal. Intersect its open neighborhood with the
two genuine tangent opens. The inverse identities, ideal invariance,
contact-coordinate identity, carrier invariance and trivial action
restrict to this intersection. Both actual glueing neighborhoods are
therefore contained in the two tangent opens and contain $x$, since the
maps fix the original cosupport.

For every actual coefficient-homogenized test prefix, the genuine
universal-property lifting construction produces a biholomorphism of
the two inverse-image neighborhoods, whole transformed coefficient-ideal
transport, whole contact-record transport, and pointwise fixation of
the actual transformed cosupport. The true contact-stage coordinate
theorem, applied inside the actual tangent opens, produces closed smooth
strict contact carriers on both neighborhoods. The transformed coefficient
ideal is coherent by the actual coherent coefficient construction and
the actual test tower; its marking is positive by the genuine factorial
coefficient marking. Thus all premises of whole remainder-record
transport have been produced internally. Apply it to the actual whole
contact-record equality for each $k$. This preserves the actual
exceptional boundary and exactly the padded historical positions.
It imposes no equality on undeleted original boundary members.
\end{proof}

\begin{theorem}[The padded datum is exactly the original boundary-kept contact remainder]
\label{hir:zero-contact-remainder-literal-history-dictionary}
\leanok
\lean{Hironaka.HironakaStepOneZeroContactRemainderGlueing.contactRemainderMarked_contactMarkedKept_emptyE}
\uses{hir:zero-contact-actual-remainder-definition,hir:drop-original-members}
At every tower prefix, restricting the contact datum of the coefficient
homogenization with empty initial boundary, removing its actual selected
components, and padding $\#K.E$ empty entries gives exactly the actual
remainder of the original boundary-kept contact datum with its first
$\#K.E$ entries replaced by empty entries.
\end{theorem}
\begin{proof}
Initially the coefficient homogenization with empty initial boundary
is the original bar datum with all original boundary entries dropped.
The actual map-and-append transformation recursion commutes with dropping
these entries: new exceptional entries are appended after the original
block, whose length persists at every prefix. The coefficient operation
and actual open restriction preserve the carrier, ideal and mark
identities; taking preimages of boundary members commutes with dropping.
The two actual contact records consequently have the same carrier,
ideal, mark and cosupport, and their boundaries differ precisely by
dropping the original block. They select precisely the same actual
contact components. Padding the dropped contact datum by the original
number of empty entries gives the defined original contact remainder
in all four fields. No exceptional entry is dropped by this padding.
\end{proof}

\begin{theorem}[Original chosen contact inputs produce the actual remainder glueing]
\label{hir:zero-contact-remainder-original-chosen-glueing}
\leanok
\lean{Hironaka.HironakaStepOneZeroContactRemainderGlueing.exists_actual_chosen_contact_remainder_glueing}
\uses{hir:zero-contact-remainder-actual-contact-lifts,hir:zero-contact-remainder-literal-history-dictionary,hir:step-one-bridge,hir:step-one-lower-domain}
Let $K$ have its actual domain-three and maximal-order properties, and
let $R$ be a canonical sequence for its literal constructed Step One
invariant. The genuine chosen contact directions at original cosupport
points $x,y$, with $x$ inside the actual chosen tangent open at $y$,
internally produce whole remainder dictionaries for the SAME original
boundary-kept contact records, with their historical block replaced by
empty entries, along every literal prefix of $R$. The same actual contact maps and
actual lifts fix the original transformed cosupport.
\end{theorem}
\begin{proof}
The actual domain gives coherent ideal, closed smooth carrier, positive
marking and well-formedness. Maximal order at the actual original
cosupport points produces tangent directions, so the literal chosen
directions are genuine and their opens contain their own points. The
stated actual incidence puts $x$ in the second chosen open as well.
Hence the preceding construction applies to these two actual directions,
without any caller-supplied contact-domain or glueing witness.

The genuine Step One canonical bridge internally makes the literal
sequence a test sequence for its actual bar datum. Deleting only its
initial historical boundary gives a true test sequence for the
coefficient-homogenized datum, whose carrier, ideal and marking have
the same literal transformation recursion. Every literal prefix inherits
this genuine test property. Apply the preceding actual lifts at each
prefix, and identify the bar and coefficient-homogenized transformed
cosupports by the genuine weak-equivalence and Step One bridge. This
gives precisely the stated joint whole record dictionaries and pointwise
fixation on the actual original cosupport. No enlarged remainder
canonicality has been used or asserted.
\end{proof}

\paragraph{Scope.}
The conclusions are actual local contact transport and entire remainder
record dictionaries. Canonicality on enlarged contact neighborhoods,
deeper zero-contact histories, descent of all zero-level thresholds
across contact patches, whole General Three analytic superlevels and
general noncompact global termination remain open.
